\chapter{File Uploads (Multer + Cloudinary)}

\section{The Upload Pipeline}

Uploading a profile avatar or a task attachment involves three separate
concerns: getting the raw file bytes off the HTTP request (Multer), storing
them somewhere durable and CDN-backed (Cloudinary), and saving the resulting
URL on a MongoDB document so the rest of the app can reference it like any
other field.

\begin{diagrambox}[Upload Pipeline]
\begin{verbatim}
React (file input) --> FormData --> POST multipart/form-data
        |
        v
   Multer middleware (parses multipart body, exposes req.file)
        |
        v
   Cloudinary (cloudinary.uploader.upload)
        |
        v
   secure_url returned
        |
        v
   Save url on User/Task document in MongoDB
\end{verbatim}
\end{diagrambox}

\subsection{Why multipart/form-data}

JSON bodies (\texttt{application/json}) can only carry text. Binary file
data is sent using the \texttt{multipart/form-data} encoding, which packs
each form field --- text or file --- into its own labeled part of the request
body. \texttt{express.json()} cannot parse this format at all, which is why
a dedicated middleware (Multer) is required specifically for routes that
accept files.

\section{Backend: Multer Configuration}

Multer can buffer files in memory or write them to disk. Since the buffer is
being immediately forwarded to Cloudinary rather than kept locally, memory
storage avoids an unnecessary disk write.

\begin{lstlisting}[language=JavaScript]
// backend/src/middleware/upload.middleware.js
const multer = require("multer");

const storage = multer.memoryStorage();

const fileFilter = (req, file, cb) => {
  const allowedMimeTypes = ["image/jpeg", "image/png", "image/webp"];
  if (allowedMimeTypes.includes(file.mimetype)) {
    cb(null, true);
  } else {
    cb(new Error("Only JPEG, PNG, or WEBP images are allowed"), false);
  }
};

const upload = multer({
  storage,
  fileFilter,
  limits: { fileSize: 2 * 1024 * 1024 }, // 2 MB max
});

module.exports = upload;
\end{lstlisting}

\subsection{Cloudinary Configuration}

\begin{lstlisting}[language=JavaScript]
// backend/src/config/cloudinary.js
const cloudinary = require("cloudinary").v2;

cloudinary.config({
  cloud_name: process.env.CLOUDINARY_CLOUD_NAME,
  api_key: process.env.CLOUDINARY_API_KEY,
  api_secret: process.env.CLOUDINARY_API_SECRET,
});

module.exports = cloudinary;
\end{lstlisting}

\begin{notebox}[NOTE]
\texttt{CLOUDINARY\_CLOUD\_NAME}, \texttt{CLOUDINARY\_API\_KEY}, and
\texttt{CLOUDINARY\_API\_SECRET} come from the Cloudinary dashboard and
belong in \texttt{.env}, never committed to source control (see
Chapter~24 on environment configuration).
\end{notebox}

\subsection{Controller: Upload and Save}

\begin{lstlisting}[language=JavaScript]
// backend/src/controllers/user.controller.js
const cloudinary = require("../config/cloudinary");
const User = require("../models/User");

const uploadAvatar = async (req, res, next) => {
  try {
    if (!req.file) {
      return res.status(400).json({ message: "No image file provided" });
    }

    // Convert the in-memory buffer to a data URI Cloudinary can accept
    const base64 = req.file.buffer.toString("base64");
    const dataUri = `data:${req.file.mimetype};base64,${base64}`;

    const result = await cloudinary.uploader.upload(dataUri, {
      folder: "task-app/avatars",
      resource_type: "image",
    });

    const user = await User.findByIdAndUpdate(
      req.user.id,
      { avatar: result.secure_url },
      { new: true }
    ).select("-password");

    res.status(200).json({ user });
  } catch (error) {
    next(error);
  }
};

module.exports = { uploadAvatar };
\end{lstlisting}

\subsection{Route Wiring}

\begin{lstlisting}[language=JavaScript]
// backend/src/routes/user.routes.js
const express = require("express");
const upload = require("../middleware/upload.middleware");
const { uploadAvatar } = require("../controllers/user.controller");
const { protect } = require("../middleware/auth.middleware");

const router = express.Router();

router.post("/me/avatar", protect, upload.single("avatar"), uploadAvatar);

module.exports = router;
\end{lstlisting}

\begin{warnbox}[WARNING]
The field name passed to \texttt{upload.single("avatar")} must exactly match
the key the frontend appends to \texttt{FormData}. A mismatch results in
\texttt{req.file} being \texttt{undefined} with no error thrown.
\end{warnbox}

\section{Frontend: Building and Sending FormData}

\begin{lstlisting}[language=JavaScript]
// frontend/src/components/AvatarUploader.jsx
import { useState } from "react";
import api from "../services/api";

function AvatarUploader({ onUploaded }) {
  const [file, setFile] = useState(null);
  const [uploading, setUploading] = useState(false);

  const handleSubmit = async (e) => {
    e.preventDefault();
    if (!file) return;

    const formData = new FormData();
    formData.append("avatar", file); // key MUST match upload.single("avatar")

    setUploading(true);
    try {
      const { data } = await api.post("/users/me/avatar", formData, {
        headers: { "Content-Type": "multipart/form-data" },
      });
      onUploaded(data.user.avatar);
    } catch (error) {
      console.error("Upload failed:", error.response?.data?.message);
    } finally {
      setUploading(false);
    }
  };

  return (
    <form onSubmit={handleSubmit}>
      <input
        type="file"
        accept="image/png, image/jpeg, image/webp"
        onChange={(e) => setFile(e.target.files[0])}
      />
      <button type="submit" disabled={!file || uploading}>
        {uploading ? "Uploading..." : "Upload Avatar"}
      </button>
    </form>
  );
}

export default AvatarUploader;
\end{lstlisting}

\begin{tipbox}[TIP]
Axios will set the correct \texttt{multipart/form-data} boundary
automatically when you pass a \texttt{FormData} instance as the body --
explicitly setting the \texttt{Content-Type} header is optional in most
setups, but harmless to include for clarity.
\end{tipbox}
